\documentclass{book}
\usepackage[utf8]{inputenc}
\usepackage[T1]{fontenc}
\usepackage{geometry}
\usepackage{graphicx}
\usepackage{float}
\usepackage{hyperref}
\usepackage[english,indonesian]{babel}

\title{Optimasi Pemrosesan Template LaTeX Menggunakan Pendekatan Berbasis AI}
\author{Academic Support Agent}
\date{\today}

\begin{document}
\maketitle

\begin{abstract}

\end{abstract}

\chapter{PENDAHULUAN}


\chapter{TINJAUAN PUSTAKA}


% Original content (commented out):
% %%%%%%%%%%%%%%%%%%%%%%%%%%%%%% 
%%% Bachelor Thesis Template
%%%%%%%%%%%%%%%%%%%%%%%%%%%%%% 

\documentclass[10pt,swedish]{book}

%%%%%%%%%%%%%%%%%%%%%%%%%%%%%% Loading packages that alter the style
\usepackage[]{graphicx}
\usepackage[]{color}
% Changing colors of text in knitr chunks
\definecolor{shadecolor}{rgb}{.97, .97, .97}
\definecolor{messagecolor}{rgb}{0, 0, 0}
\definecolor{warningcolor}{rgb}{1, 0, 1}
\definecolor{errorcolor}{rgb}{1, 0, 0}
\usepackage{alltt}
\usepackage[T1]{fontenc}
\usepackage[utf8]{inputenc}
\setcounter{secnumdepth}{3}
\setcounter{tocdepth}{3}
\setlength{\parskip}{\smallskipamount}
\setlength{\parindent}{0pt}

% Set page margins, top one is alternating page margins for book bindings, second is for electronic publishing
\usepackage[left = 3.5cm, right = 2.5cm]{geometry}
%\usepackage[left = 2.5cm, right = 2.5cm]{geometry}

% Package used for placeholder text
\usepackage{lipsum}

% Prevents LaTeX from filling out a page to the bottom
\raggedbottom

% Adding both languages, Swedish and English, so they can be used intermittently in for example abstracts.
\usepackage[swedish, english]{babel}

% All page numbers positioned at the bottom of the page
\usepackage{fancyhdr}
\fancyhf{} % clear all header and footers
\fancyfoot[C]{\thepage}
\renewcommand{\headrulewidth}{0pt} % remove the header rule
\pagestyle{fancy}

% Changes the style of chapter headings
\usepackage{titlesec}
\titleformat{\chapter}
   {\normalfont\LARGE\bfseries}{\thechapter.}{1em}{}
% Change distance between chapter header and text
\titlespacing{\chapter}{0pt}{50pt}{2\baselineskip}

% Adds table captions above the table per default
\usepackage{float}
\floatstyle{plaintop}
\restylefloat{table}

% Adds space between caption and table
\usepackage[tableposition=top]{caption}

% Adds hyperlinks to references and ToC
\usepackage{hyperref}
\hypersetup{hidelinks,linkcolor = black} % Changes the link color to black and hides the hideous red border that usually is created

% If multiple images are to be added, a folder (path) with all the images can be added here
\graphicspath{ {images/} }

% Separates the first part of the report/thesis in Roman numerals
\frontmatter

\makeatother

\IfFileExists{upquote.sty}{\usepackage{upquote}}{}
\begin{document}
\selectlanguage{swedish}

\begin{titlepage}
	\clearpage\thispagestyle{empty}
	\centering
	\vspace{2cm}

	% Rubriker
	{\large Kandidatuppsats i Statistik \par}
	\vspace{4cm}
	{\Huge \textbf{Titel på uppsatsen}} \\
	\vspace{1cm}
	{\large \textbf{Eventuell underrubrik} \par}
	\vspace{4cm}
	{\normalsize Författare1  \\ % Specifies a new line
	             Författare2  \par}
	\vspace{2cm}

    \includegraphics[scale=0.75]{liu_logga.png}

    \vspace{2cm}
    
	% Information om universitetet.
	{\normalsize Avdelningen för Statistik och maskininlärning \\ 
		Institutionen för datavetenskap \\
		Linköpings universitet \par}
	{\normalsize Vårterminen, 20XX | ISRN-nummer \par}
	\vspace{2cm}
	
	\pagebreak
	\clearpage\thispagestyle{empty}
	\raggedright
	{\tiny \ \par}
	\vspace{15cm}
	
	{\normalsize Handledare: NAMN, titel \par}
	{\normalsize Examinator: NAMN, titel \par}

\end{titlepage}

% Resets the page counter to 1
\setcounter{page}{1}

\subsection*{Sammanfattning}

Högst en sida om bakgrund, syfte och de viktigaste resultaten. Sammanfattningen
ska kunna läsas separat och ska inte innehålla några hänvisningar
till texten i övrigt. Den ska vara skriven med fullständiga meningar.
Om rapportspråket är svenska ska en sammanfattning skrivas även på
engelska. 

\cleardoublepage{}

% Changes the language of the abstract to English
\selectlanguage{english}

\subsection*{Abstract}

Abstract in English.

\cleardoublepage{}

% Changes back the language for the rest of the document to Swedish
\selectlanguage{swedish}

\subsection*{Förord}

Tack till organsationen som har tillhandahållit data, tack till\dots{}
etc.

\cleardoublepage{}

\tableofcontents{}

\cleardoublepage{}

\listoffigures

\cleardoublepage{}

\listoftables

\cleardoublepage{}

% The main body text starts here
\mainmatter

\chapter{Introduktion}

Kortfattad beskrivning av vad kapitlet innehåller.

\section{Bakgrund}

I bakgrunden ges en introduktion till ämnet. En litteraturgenomgång
av tidigare studier ges samt en motivering till problemställningen
i uppsatsen. Referenser skrivs in i texten. Exempel på hur referenser
skrivs in (manuellt) i texten (APA-stil):

Många studier har hittat negativa effekter av föräldrars socioekonomiska
bakgrund på barns utbildningsresultat och ekonomi i vuxenlivet (Dubow,
Boxe, and Huessmann, 2009; Fergusson, Horwood, \& Boden, 2008; Gibb,
Fergusson, \& Horwood, 2012; Myrberg \& Rosén, 2009). Barn från lägre
socioekonomiska grupper har dessutom fler syskon i genomsnitt (Blake,
1981; Downey, 1995; 2001). Ferrer-Wreder, Wänström och Corovic (2014)
har funnit samband mellan utbildningsresultat och \dots{} 

Med en automatisk bibliografi används \verb=\cite{}= med de enskilda källornas nycklar.

\section{Syfte}

Syfte och problemställningar med eventuella avgränsningar tas upp.

\section{Etiska och samhälleliga aspekter}

Etiska aspekter, exempelvis vad som gjorts för att enskilda observationer
inte ska kunna identifieras, tas upp här. Även samhälleliga aspekter,
såsom samhällelig nytta av studien, tas upp.

\clearpage{}

\chapter{Data}

Här ges en beskrivning av var data kommer ifrån. Här ges också en
kortfattad beskrivning av variabler och observationer i det givna
datamaterialet samt av eventuella transformationer och skapande av
nya variabler. Eventuellt kan tabeller och/eller diagram över data
presenteras. 

\clearpage{}

\chapter{Metod}

I metodkapitlet ges en beskrivning av de statistiska metoderna samt
utförande av dessa, med referenser till publikationer, som används
i uppsatsen. Tänk på att beskrivningen ska vara på en sådan nivå att
era kurskamrater kan läsa och förstå. Man kan således förvänta sig
att de vet hur man beräknar grundläggande saker så som medelvärden
och standardavvikelser (varför de inte behöver beskrivas). Kapitlet
kan delas upp i underrubriker för olika statistiska metoder.

\section{Multilevelmodeller}

Här ges en beskrivning av den (första) statistiska metoden med referenser
i texten. Här ges även numrerade formler. Exempel:

Responsen för det i:te barnet i den j:te klassen modelleras som

\begin{equation}
y_{ij}=\beta_{0j}+\beta_{1j}x_{1ij}+\varepsilon_{ij}
\end{equation}
\label{Formula Multilevel}

där:

\[
\beta_{0j}=\beta_{0}+v_{j}
\]

\[
\beta_{1j}=\beta_{1}+u_{j}
\]

Tänk på att alla termer i alla formler ska definieras samt att eventuella
antaganden ska redovisas. 

När man hänvisar till en formel anger man numret såsom: Vi kan se
i ekvation \ref{Formula Multilevel} att\dots{} 

Man kan också dela in metoden i fler underrubriker om man exempelvis
vill beskriva relevanta mått (som exempelvis Intraklasskorrelation,
Statistiska test, etc.)

\section{Den andra statistiska metoden}

\lipsum[1-3]

\subsection{Ett mått som används för denna metod}

\lipsum[3-5]

\clearpage{}

\chapter{Resultat}

I den här delen presenteras alla resultat. Resultatdelen kan delas
in i underrubriker där exempelvis beskrivande statistik redovisas
först följt av resultat från olika analyser. Alternativt kan beskrivande
statistik redovisas i Data-kapitlet.

I resultatdelen redovisas även tabeller och diagram. Alla tabeller
och diagram som redovisas måste refereras till i texten. Exempel:

I tabell \ref{Table Estimates} nedan redovisas resultaten från en
multilevelanalys av modell . Det finns en signifikant positiv effekt
av föräldrars utbildning på barns genomsnittsbetyg $(\hat{\beta}_{1}=0.5317,\,p=0.0001)$
vilket betyder att ju högre utbildning barnets föräldrar har, desto
högre genomsnittsbetyg har barnet i genomsnitt, givet att föräldrarna
har samma inkomst.

\begin{table}[H]
\begin{centering}
\begin{tabular}{|c|c|c|c|}
\hline 
Parameter &
Skattning &
Standardfel &
p-värde\tabularnewline
\hline 
\hline 
$\beta_{1}$ &
0.5317 &
0.0560 &
0.0001\tabularnewline
\hline 
$\beta_{2}$ &
0.0571 &
0.0488 &
0.2418\tabularnewline
\hline 
$\sigma_{u}^{2}$ &
0.0097 &
0.0056 &
0.0424\tabularnewline
\hline 
$\sigma_{v_{1}}^{2}$ &
... &
... &
...\tabularnewline
\hline 
$\sigma_{v_{2}}^{2}$ &
... &
... &
...\tabularnewline
\hline 
$\sigma_{v_{1}u}^{2}$ &
... &
... &
...\tabularnewline
\hline 
$\sigma_{v_{2}u}^{2}$ &
... &
... &
...\tabularnewline
\hline 
$\sigma_{v_{1}v_{2}}^{2}$ &
... &
... &
...\tabularnewline
\hline 
$\sigma_{\varepsilon}^{2}$ &
... &
... &
...\tabularnewline
\hline 
\end{tabular}
\par\end{centering}
\caption{Skattningar, standardfel och p-värden från en analys av modell }
\label{Table Estimates}
\end{table}

Innebörden/analysen av resultaten ges inte i resultatdelen, utan sparas
till diskussionsdelen.

\clearpage{}

\chapter{Diskussion/Analys}

Detta är en av de viktigaste delarna i uppsatsen. Här ingår en djupare
tolkning/innebörd av resultaten. Här ingår också kritisk granskning
av resultaten och metoderna samt en utvärdering av lämpligheten av
de använda metoderna. Här görs även kopplingar tillbaka till tidigare
studier som gåtts igenom i bakgrunden med referenser i texten. 

\clearpage{}

\chapter{Slutsatser}

Här sammanfattas slutsatserna och frågeställningarna från syftet besvaras.
Ingen ny information ska presenteras här.

\clearpage{}

\chapter*{Bibliografi}

\textit{Referenser till böcker, vetenskapliga artiklar etc. som refereras
till i uppsatsen listas här. Använd en konsekvent stil (exempelvis
APA, Harvard\dots ). Ange referenserna i bokstavsordning efter författarens
efternamn. Finns inget författarnamn kan företagsnamn användas. Allt
som refereras i uppsatsen ska finnas med i referenslistan, och allt
som finns med i referenslistan ska refereras till i uppsatsen. 
\\[\baselineskip]
Exempel:
}

Dubow, E. F., Boxer, P, \& Huessmann, L, R. (2009). Long-term effects
of parents' education on children's educational and occupational success.
Mediation by family interactions, child aggression and teenage aspirations.
Journal of Developmental Psychology, 55, 224-249. 

Fergusson, E. M., Horwood, L. J., \& Boden, J. M. (2008). The transmission
of social inequality: Examination of the linkages between family socioeconomic
status in childhood and educational achievement in young adulthood.
Research in Social Stratification and Mobility, 26, 277-295. 

Montgomery, D.C. (1991). Design and Analysis of Experiments, Wiley,
New York. 
\\[\baselineskip]

\textit{De första två referenserna är artiklar i tidskrifter. Där anges efternamn,
initialer, årtal i parentes, artikeltitel, tidsskriftstitel i kursivt,
volym i kursivt, samt sidnummer. Den tredje referensen är en bok.
Där anges efternamn, initialer, årtal i parantes, boktitel kursivt,
förlag och ort. 
\\[\baselineskip]
Referenserna behöver inte delas upp i böcker och artiklar. Försök
att använda publikationer så mycket som möjligt, dvs publicerade böcker,
artiklar, utredningar etc. Hittar ni något på en internetsida \textendash{}
undersök om det finns som en publikation någonstans.
\\[\baselineskip]
Ni kan också använda er av Mendeley som ett hjälpmedel för att organisera
och hänvisa till era källor. Exportera då databasen till .bibtex-format
och använd de inbyggda citationsverktygen som finns i LyX/LaTeX.}

% Adds an appendix that uses capital Roman numerals as page numbering
\appendix
\backmatter
\pagenumbering{Roman}

\chapter*{Bilaga}

I bilagorna kan man lägga detaljerade tabeller som är för stora/detaljerade
för att ha i den löpande texten, programkod etc. Bilagor kan numreras
A, B, C, etc.
\end{document}


\end{document}